% !PRINTED 629
\weberpage{629}
\gesperrt{4. Aus 2. und 3. folgt, wenn $\zeta$ von Null verschieden ist:}
\[
N(\zeta)\,N\left(\frac{1}{\zeta}\right) = 1,
\]
\gesperrt{und daraus für irgend zwei $\zeta, \zeta'$:}
\weq{7}{N\left(\frac{\zeta'}{\zeta}\right) = \frac{N(\zeta')}{N(\zeta)}.}

\gesperrt{5. Ist $t$ eine unbestimmte oder variable Größe, und $\varphi(t)$ die Funktion, die nach (3) für $t = \zeta$ zu Null wird, so ist}
\weq{8}{\varphi(t) = N(t - \zeta).}

Das ergibt sich aus (1) und (2). Denn ersetzt man in (1) $\zeta$ durch $\zeta - t$, so ist dies gleichbedeutend damit, daß man $y_{i,i}$ durch $y_{i,i} - t$ ersetzt und die übrigen $y_{i,k}$ ungeändert läßt.

Zerlegen wir die Funktion $\varphi(t)$, die eine rationale Funktion von $t$ und $z$ ist, in ein Produkt von irreduziblen Faktoren $\varphi_1(t), \varphi_2(t), \dots$ so muß einer dieser Faktoren für $t = \zeta$ verschwinden. Sei dies $\varphi_1(t)$, und
\weq{9}{\varphi_1(t) = t^e + b_1' t^{e-1} + \cdots + b_{e-1}' t + b_e'}
sei vom Grade $e$. Es ist dann $\varphi_1(\zeta) = 0$ die Gleichung niedrigsten Grades, der $\zeta$ genügt, und aus $\zeta$ leitet man einen Körper $\Omega_1$ ab, der in $\Omega$ enthalten ist und den Grad $e$ hat. Jede Zahl $\eta$ des Körpers $\Omega_1$ kann, und zwar auf eine Art, in der Form dargestellt werden:
\weq{10}{\eta = x_0 + x_1 \zeta + \cdots + x_{e-1} \zeta^{e-1}.}

Sei ferner
\weq{11}{\theta^f + \eta_1 \theta^{f-1} + \cdots + \eta_{f-1} \theta^{f-1} + \eta_f = 0}
% [sic] printed \eta_{f-1} \theta^{f-1} as the penultimate term, the powers of \theta descend so this term should be \eta_{f-1} \theta
die Gleichung niedrigsten Grades mit Koeffizienten in $\Omega_1$, der $\theta$ genügt, also $f$ der Relativgrad von $\Omega$ in bezug auf $\Omega_1$.

Durch die $ef$ Funktionen
\weq{12}{\zeta^h \theta^k \qquad \begin{array}{l} h = 0, 1, \dots, e-1 \\ k = 0, 1, \dots, f-1 \end{array}}
kann jede Funktion des Körpers linear ausgedrückt werden, und zwischen ihnen besteht keine lineare Gleichung mit rational von $z$ abhängigen Koeffizienten. Denn sonst würde $\theta$ oder $\zeta$ aus einer Gleichung von niedrigerem Grade als $f$ oder $e$ entstehen. Dem-
%CONT>>
